\documentclass[11pt]{article}
\usepackage[margin=1in]{geometry}
\usepackage{amsmath,amssymb,amsthm}
\usepackage{algorithm}
\usepackage{algpseudocode}
\usepackage{booktabs}
\usepackage{hyperref}

\newtheorem{definition}{Definition}
\newtheorem{theorem}{Theorem}
\newtheorem{lemma}[theorem]{Lemma}
\newtheorem{property}{Property}

\title{\textbf{Mathematical Specification, Coordinate Preconditioning, and Regret Bounds of the AdaGrad Adaptive Optimizer (ALGO-NN-33)}}
\author{Team Scaibu \\ \small Email: \texttt{jaydeep.9664920749@gmail.com} \\ \small Architecture and Core Engine Team}
\date{\today}

\begin{document}
\maketitle

\begin{abstract}
AdaGrad (Adaptive Gradient Algorithm) automatically scales coordinate learning rates inversely proportional to the square root of historical squared gradient sums $G_{t, i} = \sum_{\tau=1}^t g_{\tau, i}^2$. Infrequently updated features receive larger relative learning rates, while frequent coordinates are damped. We formalize the diagonal proximal metrics, prove the $\mathcal{O}(\sqrt{T})$ sublinear regret bound, provide line-by-line pseudocode matching the reference implementation, and derive complexity.
\end{abstract}

\section{Notation and Mathematical Preliminaries}
\label{sec:notation}

Let $\theta \in \mathbb{R}^P$ denote parameter coordinates, $\mathbf{g} \in \mathbb{R}^P$ parameter gradients, and $\mathbf{G} \in \mathbb{R}_{\ge 0}^P$ the accumulated squared gradient state buffer. Let $\eta \in \mathbb{R}_{>0}$ denote the global learning rate and $\epsilon \in \mathbb{R}_{>0}$ the stability denominator.

\subsection{Coordinate Adaptation Mechanics}
\paragraph{Intuitive Meaning:} When training NLP or recommendation models with huge embedding tables, rare token IDs get very few gradient updates, while common words get millions of updates. AdaGrad keeps a separate speedometer for each parameter, giving big steps to rare coordinates and small steps to common ones.

\paragraph{Formal Mathematical Definition:}
\begin{equation}
\label{eq:adagrad_update}
G_{t, i} = G_{t-1, i} + g_{t, i}^2, \quad \theta_{t, i} = \theta_{t-1, i} - \frac{\eta}{\sqrt{G_{t, i}} + \epsilon} g_{t, i}.
\end{equation}

\paragraph{Worked Numerical Example:}
Let $\theta = [0.0], G = [0.0], g = [2.0], \eta = 0.1, \epsilon = 10^{-8}$:
$G_{\mathrm{next}} = 0.0 + 2.0^2 = 4.0$.
$\Delta\theta = \frac{0.1}{\sqrt{4.0} + 10^{-8}} \times 2.0 = \frac{0.1}{2.0} \times 2.0 = 0.1$.
$\theta_{\mathrm{next}} = 0.0 - 0.1 = -0.1$.

\subsection{Summary of Symbols}
\begin{itemize}
    \item $P \in \mathbb{N}$: Parameter dimension.
    \item $\boldsymbol{\theta} \in \mathbb{R}^P$: Parameter vector.
    \item $\mathbf{g} \in \mathbb{R}^P$: Gradient vector.
    \item $\mathbf{G} \in \mathbb{R}^P$: Accumulated squared gradient state buffer.
    \item $\eta \in \mathbb{R}_{>0}$: Global learning rate.
    \item $\epsilon \in \mathbb{R}_{>0}$: Epsilon denominator ($10^{-8}$).
\end{itemize}

\section{Problem Definition and Core Concepts}
\label{sec:problem_definition}

\subsection{Diagonal Metric Tensor Interpretation}
AdaGrad is an instance of online mirror descent where the local Riemannian metric tensor is $\mathbf{H}_t = \operatorname{diag}(\sqrt{\mathbf{G}_t})$. The effective step size along coordinate $i$ decays monotonically as $\eta_{t, i} \propto \frac{1}{\sqrt{t}}$.

\section{Algorithmic Specification}
\label{sec:algorithm}

Algorithm~\ref{alg:adagrad} specifies the coordinate update matching \texttt{impl.py}.

\begin{algorithm}[H]
\caption{AdaGrad Adaptive Optimizer Step Evaluation}
\label{alg:adagrad}
\begin{algorithmic}[1]
\Require Parameters $\boldsymbol{\theta} \in \mathbb{R}^P$, gradients $\mathbf{g} \in \mathbb{R}^P$, accumulator $\mathbf{G} \in \mathbb{R}^P$, learning rate $\eta > 0$, epsilon $\epsilon > 0$.
\Ensure Updated parameters $\boldsymbol{\theta}_{\mathrm{next}} \in \mathbb{R}^P$, updated accumulator $\mathbf{G}_{\mathrm{next}} \in \mathbb{R}^P$.
\State Initialize $\boldsymbol{\theta}_{\mathrm{next}} \gets [], \quad \mathbf{G}_{\mathrm{next}} \gets []$
\For{$i = 1$ to $P$}
    \State $G_{\mathrm{new}} \gets G_i + g_i^2$
    \State $step \gets \frac{\eta}{\sqrt{G_{\mathrm{new}}} + \epsilon} \cdot g_i$
    \State $\theta_{\mathrm{new}} \gets \theta_i - step$
    \State Append $\theta_{\mathrm{new}}$ to $\boldsymbol{\theta}_{\mathrm{next}}$, append $G_{\mathrm{new}}$ to $\mathbf{G}_{\mathrm{next}}$
\EndFor
\State \Return $\boldsymbol{\theta}_{\mathrm{next}}, \quad \mathbf{G}_{\mathrm{next}}$
\end{algorithmic}
\end{algorithm}

\section{Theoretical Guarantees and Regret Bounds}
\label{sec:guarantees}

\begin{theorem}[AdaGrad $\mathcal{O}(\sqrt{T})$ Regret Bound]
\label{thm:adagrad_regret}
\textbf{Assumptions:} Let loss functions $f_t$ be convex with bounded subgradients $\|g_t\|_\infty \le G_{\infty}$ and bounded parameter diameter $\|\theta - \theta^*\|_2 \le D$.
\textbf{Guarantee:} AdaGrad achieves total regret $R(T) = \sum_{t=1}^T f_t(\theta_t) - \min_{\theta^*} \sum_{t=1}^T f_t(\theta^*) \le \sqrt{2} D \sum_{i=1}^P \sqrt{\sum_{t=1}^T g_{t, i}^2} = \mathcal{O}(\sqrt{T})$.
\textbf{Proof:}
Using the standard potential function $\Phi_t(\theta) = \frac{1}{2} (\theta - \theta^*)^\top \operatorname{diag}(\sqrt{\mathbf{G}_t}) (\theta - \theta^*)$, summing telescoping terms over $t$ and bounding coordinate sums using the integral inequality $\sum_{t=1}^T \frac{g_{t, i}^2}{\sqrt{G_{t, i}}} \le 2\sqrt{G_{T, i}}$ yields the sublinear bound $\mathcal{O}(\sqrt{T})$.
\textbf{Limitations:} Monotonic accumulation $G_{t, i} \ge G_{t-1, i}$ causes the effective step size to continuously decay toward zero, stopping learning prematurely on long non-convex training runs (resolved by RMSProp and Adam).
\end{theorem}

\section{Complexity Analysis}
\label{sec:complexity}

\begin{itemize}
    \item \textbf{Time Complexity:} $\mathcal{O}(P)$ operations per step.
    \item \textbf{Space Complexity:} $\mathcal{O}(P)$ memory for the squared gradient accumulator buffer $\mathbf{G}$.
\end{itemize}

\section{Worked Numerical Example}
\label{sec:worked_example}

Let $\theta = [1.0], G = [1.0], g = [3.0], \eta = 0.5, \epsilon = 0.0$:
\begin{enumerate}
    \item $G_{\mathrm{new}} = 1.0 + 3.0^2 = 1.0 + 9.0 = 10.0$.
    \item $\sqrt{G_{\mathrm{new}}} = \sqrt{10.0} \approx 3.162278$.
    \item $\Delta\theta = \frac{0.5}{3.162278} \times 3.0 = 0.158114 \times 3.0 = 0.474342$.
    \item $\theta_{\mathrm{next}} = 1.0 - 0.474342 = 0.525658$.
\end{enumerate}

\section{Numerical Considerations and Edge Cases}
\label{sec:numerical}

\begin{itemize}
    \item \textbf{Zero Gradient Step:} When $g_i = 0$, $G_{\mathrm{new}} = G_i$ and $\Delta\theta = 0$, guaranteeing numerical stability without drift.
\end{itemize}

\begin{thebibliography}{9}
\bibitem{duchi2011} J.~Duchi, E.~Hazan, and Y.~Singer, ``Adaptive Subgradient Methods for Online Learning and Stochastic Optimization,'' \emph{JMLR}, vol.~12, pp.~2121--2159, 2011.
\end{thebibliography}

\end{document}
